\chapter{Conclusion}\label{ch:conclusion}

This thesis was based on how far the mapping of security requirements between
\standardA{} and \standardB{} can be automated with current NLP and LLM
techniques. A five-stage pipeline was set for this thesisi project that extracts 141 normative
provisions from the two standards, compare and maps them with eleven different methods and 
evaluate their performance. The methods covers lexical similarity, sentence-embedding, contextual-embedding, cross-encoder and hosted API approaches. The evaluation was based on  a 200-pair reference set assembled
from three rounds of LLM-assisted annotation with totalling 1{,}027 judged pairs. 
Every label drafted by one frontier model, after that the annotations were analyzed and reviewed 
when necessary, adjudicated by the authors wherever the two models diverged.


The results shows that automated mapping can help to identify whether two provisions are related,
but it is not yet reliable enough to determine the exact type of relationship between them. 
In other words, finding potentially provisions is partly feasible, while assigning a meaningful relationship between them still requires human judgement.


The detection result turns on calibration. The reference set was built to be an equal number of related and unrelated pairs, so the classifier that calls every pair related scores $F_1 = 0.667$ on it, and 
that figure is the bar every method has to clear. When their default thresholds are used 
four of the eleven methods does not exceed this baseline. SecureBERT, CySecBERT, Gemini embeddings and, to within one pair, BERT mark \emph{every} annotated pair as related and reproduce the trivial score exactly and performs no better 
than the simple baseline.

Choosing each threshold by cross-validation improves the results.
Gemini embeddings reach F1 0.782 at precision 0.774, a margin of $+0.116$ [0.050, 0.188] over the baseline; BGE and CySecBERT improves its baseline by smaller margins whose intervals also exclude zero. As per the given though at a
reference set measuring $\kappa = 0.51$, these smaller differences should not be interpreted as strong evidence that one method is better than another. The major practical finding is that calibration can help to prevent optimistic benchmarks; a model from behaving like an almost constant classifier. However, calibration alone does not turn the relatively small improvements into highly reliable automation mapping.

Determining how two provisions relate is not solved by any of the eleven evaluated methods performs better than the majority-class baseline, which achieves a typed accuracy of 0.505. 
This suggests that a single similarity score is not sufficient to distinguish between different semantic relationships. The reason is the framing rather than the representations: a control trained on the 827 judged pairs held out of the reference set, using only the scores those same eleven methods produce,
reaches 0.580 accuracy and macro-F1 0.321 against a trivial 0.500 and 0.133 for the majority-class baseline. This indicates that the similarity scores contain some information about relationship types, but that the fixed similarity thresholds used by the individual methods do not capture this information effectively.

What that control does not do is resolve the overlap/complementarity boundary,
where it erorrs as the annotators do, so typed mapping remains well short of what
compliance documentation requires. No similarity-based method can represent
directional subsumption even in principle, and Unlike the other methods, the cross-encoder was calculated in both directions which will able to identify subsumption relationships. It produced 123 subsumption labels, whereas the other ten methods produced none. However, its recall was too low for the method to be practically useful, and its computational cost was approximately sixty times higher. This proves that more complex models can capture information that scalar similarity methods miss, but it shows that with higher model complexity does not automatically provide sufficiently reliable results.

Two main technical mechanisms are observed across the experiments. First the
general-purpose embedding spaces suffer from anisotropy, which drives BERT-family models to assign high similarity to almost all pairs and secondly, the mismatch between scalar similarity and a typed relationship taxonomy. Cybersecurity-specific pretraining did not help in either of the two independent attempts tested (SecureBERT, CySecBERT), while task-specific training for sentence comparison including Sentence-BERT, MPNet, and BGE, generally performed better. Within the models evaluated in this thesis, this suggests that how a model was trained for the comparison task was more important than whether it was specifically trained on cybersecurity text.

A third finding came from the annotation rather than the models. Equivalence and
subsumption are close to absent from the collective data. Only 13 pairs among the 1{,}027
judged, and none at all among the 300 highest-ranked unjudged pairs annotated
specifically to identify more examples. A device-level standard and a lifecycle-level standard
produce overlapping and complementary obligations in quantity, and identical or
strictly-containing ones almost never. For this particular pair of standards, the six-label relationship taxonomy effectively behaves as a three-label taxonomy.

Overall, the results indicates that automated compliance mapping should currently used as decision
support rather than humans assign the relationship labels or judgement. Even that reduced role is valuable to a compliance team, since the ranking step compresses 4{,}968 candidate pairs into a short review list in under
a minute of computation which can provide substantial practical value to compliance teams.

\subsection*{Critical discussion}

Three choices shaped this project and deserve scrutiny. First, the reference set
was annotated with LLM assistance; this was a pragmatic necessity at 15 credits,
it is disclosed openly, and the design (two annotation models from separate
developers, neither sharing a family with the evaluated ones, a written codebook
fixed in advance, author adjudication of every model disagreement, and a blind
comparison against the author-reviewed pilot) bounds the risk. What it does not remove is
the ceiling that comparison reveals: $\kappa = 0.51$ on the detection decision
means the reference is good enough to rank methods against a common target but
not to certify an absolute accuracy. A fully expert-built reference would be
stronger, and the survey in Appendix~\ref{app:survey} is the direct path to one.

Second, and most consequentially, the reference set was built to a positive rate
of one half. That choice succeeded in what it was for: the degeneracy of four
methods is visible directly in Table~\ref{tab:detection} instead of being hidden
behind a high number. One half is nonetheless a design parameter and the corpus is
far sparser, so every performance figure in this thesis stands as an upper bound
on what the same operating point achieves across all 4{,}968 candidates
(Section~\ref{sec:refset}). The size of that gap cannot be measured from this
design. The coverage analysis of Section~\ref{sec:coverage} should be read
accordingly, as the shape of the predicted mapping and not its accuracy at corpus
scale.

Third, reporting two operating points per method rather than one doubles the
tables but is the only way to separate what a model can represent from where its
threshold happens to sit.

\subsection*{Generalisation of the results}

The specific numbers belongs to this standard pair, reference set and these
model versions. The structural findings travel further. An isotropy of mean pooled
contextual embeddings is a documented model property. The symmetric based method cant not express asymmetric relationships when one particular requirement is more detailed and specific than another requirement. The limitation can be easily easily visible wherever two related security
frameworks are mapped. The pipeline itself is standard agnostic, and any pair of
documents with identifiable normative clauses can be replaced.

\subsection*{Future work}\label{sec:futurework}

Four directions follows naturally, and the first is a direct repair of this design's principal limitation.

First, model measures the performance at the corpus base rate. Every figure here is
conditioned on a reference set built to be half positive, and the gap between
that and the sparse reality of 4{,}968 candidate pairs is currently unquantified. Due to this difference, the current outcomes may not cover whole real world performance. The better approaches might be
drawing a sample with known inclusion probabilities over the full space and
estimating precision, recall and prevalence from it would convert the upper
bounds reported here into estimates, and would put a number on how far the
calibrated thresholds of Table~\ref{tab:calibrated} drift when applied where
positives are rare. This is the single change that would most improve the
strength of the claims.

Second, execute the expert validation survey designed in
Appendix~\ref{app:survey} to put the ground truth on a multi-rater footing and to
calibrate the human agreement ceiling.

Third, train the typed task properly. The control in
Section~\ref{sec:supervised} shows that supervision over the stored scores alone
lifts typed accuracy above every zero-shot method, which makes end-to-end training
the obvious extension rather than a speculative one: a cross-encoder fine-tuned on
requirement pairs, rather than borrowed from general natural language inference,
would train the representation and the decision together, and entailment scoring
is in any case the only approach here that recovers subsumption direction. The
scarcity finding tempers this: with thirteen rare-class pairs in the whole corpus,
the signal for those two classes would have to come from other standard pairs.

Fourth, revisit the prompt given to the LLM classifier. Its overuse of
complementarity tracks the wording of the label definition rather than the
difficulty of the pairs, since ``related but different concerns that work
together'' admits almost any two security requirements. The ceiling measured here
may therefore be the definition's rather than the model's.

Beyond these, ranking metrics such as precision at $k$ and mean average precision
would decouple the evaluation from the positive rate altogether, since they
measure the shortlisting behaviour the pipeline is actually fit for; the
machinery is implemented but the judged fraction of each ranking is still so thin
that precision among judged candidates is 1.000 for every method, meaning the
numbers describe judging coverage rather than ranking quality.
